\renewcommand{\thechapter}{ก}
\chapter{สูตรคณิตศาสตร์และตารางค่าคงที่ฟิสิกส์เกษตร}
\label{chap:appendixA}

\section{สรุปสมการหลักการแรกในฟิสิกส์เกษตรดิจิทัล}

\subsection{สมการไฟฟ้าเคมีและอุณหพลศาสตร์ของดิน}
\begin{enumerate}
    \item \textbf{สมการเนิร์นสต์ (Nernst Equation)} สำหรับการวัดค่าความเป็นกรด-ด่าง
    \begin{equation}
    E = E^0 - \frac{2.303 R T}{F} \text{pH}
    \end{equation}
    \item \textbf{สมการความดันไอน้ำอิ่มตัวของเทเทนส์ (Tetens Equation)}
    \begin{equation}
    e_s(T) = 0.61078 \exp\left( \frac{17.27 T}{T + 237.3} \right) \quad (\text{kPa})
    \end{equation}
    \item \textbf{สมการแรงดึงระเหยน้ำของอากาศ (VPD)}
    \begin{equation}
    \text{VPD} = e_s(T) \times \left( 1 - \frac{\text{RH}}{100} \right) \quad (\text{kPa})
    \end{equation}
\end{enumerate}

\subsection{สมการชลศาสตร์ท่อและการคายระเหยน้ำของพืช}
\begin{enumerate}
    \item \textbf{การคายระเหยน้ำของพืช ($ET_c$)}
    \begin{equation}
    ET_c = ET_0 \times K_c \quad (\text{mm/day})
    \end{equation}
    \item \textbf{ปริมาตรน้ำที่ต้องการต่อต้น ($V_{\text{water}}$)}
    \begin{equation}
    V_{\text{water}} = \frac{A_{\text{canopy}} \times ET_c}{\text{Eff}} \quad (\text{Liters/tree/day})
    \end{equation}
    \item \textbf{สมการการสูญเสียแรงดันในท่อของเฮเซน-วิลเลียมส์ (Hazen-Williams)}
    \begin{equation}
    h_f = 10.67 \times \frac{L \times Q^{1.852}}{C^{1.852} \times D^{4.87}} \quad (\text{meters})
    \end{equation}
    \item \textbf{เฮดรวมของระบบชลศาสตร์ ($TDH$)}
    \begin{equation}
    TDH = H_s + h_f + h_m + H_{op} \quad (\text{meters})
    \end{equation}
\end{enumerate}

\subsection{สมการอุณหพลศาสตร์การละลายของก๊าซและเครื่องกลหมุน}
\begin{enumerate}
    \item \textbf{สมการการดับการเรืองแสงสเติร์น-โวลเมอร์ (Stern-Volmer)}
    \begin{equation}
    \frac{\tau_0}{\tau} = 1 + K_{SV} [\text{O}_2]
    \end{equation}
    \item \textbf{การถ่ายเทมวลสารของออกซิเจนละลายน้ำ}
    \begin{equation}
    \frac{dC}{dt} = K_L a \left( C^*_s - C \right)
    \end{equation}
    \item \textbf{กฎความคล้ายคลึงของเครื่องกลหมุนสำหรับการประหยัดพลังงาน (Affinity Laws)}
    \begin{equation}
    \frac{P_2}{P_1} = \left( \frac{N_2}{N_1} \right)^3
    \end{equation}
\end{enumerate}
\clearpage
